% !TEX TS-program = xelatex
% !TEX encoding = UTF-8 Unicode

\documentclass{resume}

\usepackage{zh_CN-Adobefonts_external}
\usepackage{linespacing_fix}
\usepackage{cite}
\usepackage{graphicx}
\usepackage{hyperref}

\begin{document}
\pagenumbering{gobble}

\name{Zhongtao Rao}

\contactInfo
{raozhongtao@gmail.com}
{https://raozhongtao.github.io/}


% ============================================================
\section{Research Interests}

Large Recommendation Models, Recommendation Model Scaling,
User Behavior and Sequence Modeling, and Representation Learning.

% ============================================================
\section{Education}

\datedsubsection{
\textbf{Hong Kong University of Science and Technology (Guangzhou)}

$\mid$ M.Phil. in Data Science and Analytics
}{Aug. 2024 -- Jan. 2027 (Expected)}

\textbf{Research Advisor}: Prof. Nan Tang

\vspace{0.3em}

\datedsubsection{
\textbf{South China University of Technology}

$\mid$ B.Sc. in Information and Computing Science
}{Sep. 2019 -- Jun. 2023}

% ============================================================
\section{Publications}

\begin{itemize}

\item
\textbf{Zhongtao Rao}, Peilin Zhou, Dading Chong, Zihan Chen,
Shoujin Wang, Nan Tang.
\textit{Structural and Disentangled Adaptation of Large Vision Language Models
for Multimodal Recommendation}.
In Proceedings of the 49th International ACM SIGIR Conference on Research and
Development in Information Retrieval (\textbf{SIGIR 2026}).
\href{https://arxiv.org/abs/2512.06883}{[arXiv]}

\item
Peilin Zhou, Chao Liu, Jing Ren, Xinfeng Zhou, Yueqi Xie, Meng Cao,
\textbf{Zhongtao Rao}, You-Liang Huang, Dading Chong, Junling Liu,
Jae Boum Kim, Shoujin Wang, Raymond Chi-Wing Wong, Sunghun Kim.
\textit{When Large Vision Language Models Meet Multimodal Sequential
Recommendation: An Empirical Study}.
In Proceedings of the ACM Web Conference 2025
(\textbf{WWW 2025, Oral}).

\end{itemize}

% ============================================================
\section{Industry Research Experience}

% ---------------- Microsoft ----------------
\datedsubsection{
\textbf{Microsoft China $\mid$ Microsoft AI $\mid$ Applied Scientist Intern}
}{Jul. 2026 -- Present}

\textbf{Project: HyFormer-based Large Recommendation Model for News Ranking}

\begin{itemize}

\item
Investigated the scaling bottleneck of a production-scale news ranking model,
where simply widening conventional MLP-based architectures provided limited
returns due to insufficient modeling of long user behavior sequences and
interactions among heterogeneous features.

\item
Upgraded the ranking model to the \textbf{HyFormer} paradigm:
modeled user behavior sequences with full attention, introduced global tokens
with cross-attention for sequence aggregation, and jointly tokenized sparse and
dense non-sequential features to enable token-level interactions between
behavioral interests and contextual signals.

\item
Integrated sequence representations, non-sequential tokens, and dense features
into a unified interaction architecture with RankMixer and MaskNet, enabling
richer high-order feature interactions beyond conventional feature
concatenation and MLP scaling.

\item
Scaled the model from approximately \textbf{64M to 95M parameters} and trained
from cold start on approximately \textbf{4 billion samples over seven days},
achieving \textbf{+0.09\% AUC Like}, \textbf{+0.185\% RecirDwell}, and
\textbf{+0.057\% Click} in offline evaluation.

\end{itemize}

\vspace{-0.3em}
\noindent\rule{\textwidth}{0.2pt}

% ---------------- Baidu ----------------
\datedsubsection{
\textbf{Baidu $\mid$ Recommendation Strategy Department $\mid$ Algorithm Intern}
}{Nov. 2025 -- Mar. 2026}

\begin{itemize}

\item
Post-trained \textbf{Qwen3-VL-3B} with behavior-derived item-item supervision
and in-batch contrastive learning to align multimodal content representations
with user interaction signals for cold-start recommendation.

\item
Compressed 2048-dimensional representations to 100 dimensions and deployed
them in an online I2I retrieval pipeline, improving cold-start first-hop CTR
by \textbf{4.97\%}.

\end{itemize}

\vspace{-0.3em}
\noindent\rule{\textwidth}{0.2pt}

% ---------------- ByteDance Douyin ----------------
\datedsubsection{
\textbf{ByteDance $\mid$ Douyin Live $\mid$ Recommendation Algorithm Intern}
}{Apr. 2026 -- Jul. 2026}

\begin{itemize}

\item
Designed a \textbf{PID-based adaptive traffic allocation mechanism} for
high-potential long-tail creators across retrieval and ranking stages.

\item
Combined posterior creator-quality signals with real-time exposure statistics;
online experiments improved creator activity and engagement metrics.

\end{itemize}

% ============================================================
\section{Industry Experience}

\datedsubsection{
\textbf{ByteDance $\mid$ Feishu/Lark Mail $\mid$ Software Engineer}
}{Jul. 2023 -- Jul. 2024}

Developed and maintained cross-platform Android and iOS infrastructure for
the email reading module.

\end{document}

